%% ch07-log-format.tex — Log Format Reference
%% JA4proxy Technical Reference Manual

\chapter{Log Format Reference}
\index{log format}
\label{ch:logformat}

\begin{chapterquote}
\ja\ emits structured JSON logs to stdout. Every connection produces exactly one
connection log entry at the action decision point. Additional entries are emitted
for config reloads, errors, and policy events.
\end{chapterquote}

\newthought{JSON is the primary log format.} The human-readable text format is a
secondary view, derived from the JSON, suitable for interactive terminal use.
All SIEM integrations, log aggregation (Loki), and alerting pipelines should use
the JSON format.

\section{JSON Log Schema}
\index{log format!JSON}

\subsection{Complete Field Reference}

\begin{fullwidth}
\begin{table}[h]
\caption{JSON log field reference — connection event}
\label{tab:log-fields}
\begin{tabular}{@{}lllp{5.5cm}@{}}
\toprule
\textbf{Field} & \textbf{Type} & \textbf{Always present} & \textbf{Description} \\
\midrule
\code{timestamp} & string (ISO 8601) & Yes & UTC timestamp with milliseconds: \code{2026-04-04T12:34:56.789Z} \\
\code{event} & string & Yes & Event type: \code{connection}, \code{config\_reload}, \code{error}, \code{bypass\_disabled} \\
\code{action} & string & Yes & Final action: \code{allow}, \code{block}, \code{tarpit}, \code{ban}, \code{flag}, \code{rate\_limit} \\
\code{client\_ip} & string & Yes & Canonical client IP (from PROXY protocol). Full IPv4 or IPv6 string. \\
\code{country} & string & No & ISO 3166-1 alpha-2 country code. Absent if GeoIP lookup failed. \\
\code{ja4} & string & No & JA4 fingerprint string. Absent if ClientHello could not be parsed. \\
\code{ja4x} & string & No & JA4X extended fingerprint. Absent if not computed. \\
\code{ja4t} & string & No & JA4T TCP fingerprint. Absent if not computed. \\
\code{fingerprint\_name} & string & No & Human-readable label for JA4, from \configkey{security.fingerprint\_labels} \\
\code{tls\_version} & string & No & TLS version string: \code{TLS 1.3}, \code{TLS 1.2}, etc. \\
\code{sni} & string & No & Server Name Indication from ClientHello. Absent if not present in handshake. \\
\code{alpn} & string & No & First ALPN value: \code{h2}, \code{h1}, etc. \\
\code{risk\_score} & integer & No & Computed risk score (0–100). Absent for bypassed connections. \\
\code{dial} & integer & Yes & Current dial setting at decision time. \\
\code{bypass} & string & No & Bypass type if connection was bypassed: \code{alpn\_browser}, \code{ja4\_whitelist}, etc. \\
\code{reason} & string & Yes & Human-readable reason for the action. \\
\code{duration\_ms} & float & Yes & Pipeline execution time in milliseconds. \\
\code{signals} & object & No & Per-signal scores as a JSON object. Absent for bypassed connections. \\
\bottomrule
\end{tabular}
\end{table}
\end{fullwidth}

\subsection{The signals Object}
\index{log format!signals object}

When present, the \code{signals} field contains the individual score contribution
from each signal:

\begin{lstlisting}[language=json]
"signals": {
  "tls":       0,
  "sni":       0,
  "tcp":       0,
  "asn":       15,
  "dns":       0,
  "beacon":    0,
  "abuseipdb": 0,
  "rdap":      0
}
\end{lstlisting}

Each value is the raw signal score (0–100) before confidence weighting. A score of 0
means either the signal found nothing suspicious, or the external service was
unavailable (fail open — both produce 0).

\section{Example Log Entries}
\index{log format!examples}

\subsection{Allowed Browser Connection (ALPN bypass)}
\index{log format!allow example}

\begin{lstlisting}[language=json]
{
  "timestamp": "2026-04-04T12:34:56.789Z",
  "event": "connection",
  "action": "allow",
  "client_ip": "172.19.0.10",
  "country": "IE",
  "ja4": "t13d1113h2_8daaf6152771_02713d6af862",
  "fingerprint_name": "Browser (TLS 1.3)",
  "tls_version": "TLS 1.3",
  "sni": "example.com",
  "alpn": "h2",
  "dial": 0,
  "bypass": "alpn_browser",
  "reason": "Browser ALPN bypass",
  "duration_ms": 0.8
}
\end{lstlisting}

Note: \code{risk\_score} and \code{signals} are absent because the connection was
bypassed before scoring.

\subsection{Blocked Connection (Spamhaus DROP match)}
\index{log format!block example}
\index{Spamhaus!log entry}

\begin{lstlisting}[language=json]
{
  "timestamp": "2026-04-04T12:35:10.421Z",
  "event": "connection",
  "action": "block",
  "client_ip": "185.220.101.1",
  "country": "DE",
  "ja4": "t13d1900h2_8daaf6152771_97f8aa674fd9",
  "fingerprint_name": "Tool/Bot (TLS 1.3)",
  "tls_version": "TLS 1.3",
  "sni": null,
  "dial": 75,
  "bypass": "spamhaus",
  "reason": "Spamhaus DROP match: 185.220.96.0/21",
  "duration_ms": 0.3
}
\end{lstlisting}

\subsection{Scored and Tarpitted Connection}
\index{log format!tarpit example}

\begin{lstlisting}[language=json]
{
  "timestamp": "2026-04-04T12:36:42.100Z",
  "event": "connection",
  "action": "tarpit",
  "client_ip": "45.142.212.100",
  "country": "RU",
  "ja4": "t13d190900_9dc949149365_97f8aa674fd9",
  "ja4x": "a5c0d0f3d1b2_0e8a7c6b5d4e",
  "ja4t": "8192_2_1460_nWS",
  "fingerprint_name": "Sliver C2",
  "tls_version": "TLS 1.3",
  "sni": "cdn-185.example.com",
  "risk_score": 87,
  "dial": 75,
  "reason": "Score 87 exceeds tarpit threshold",
  "duration_ms": 2.4,
  "signals": {
    "tls":       0,
    "sni":       40,
    "tcp":       25,
    "asn":       65,
    "dns":       35,
    "beacon":    50,
    "abuseipdb": 80,
    "rdap":      60
  }
}
\end{lstlisting}

\subsection{Banned Connection (Subsequent Request)}
\index{log format!ban example}

\begin{lstlisting}[language=json]
{
  "timestamp": "2026-04-04T12:37:05.002Z",
  "event": "connection",
  "action": "block",
  "client_ip": "45.142.212.100",
  "country": "RU",
  "dial": 75,
  "bypass": "ban_cache",
  "reason": "Banned for 604800s (escalation level 3)",
  "duration_ms": 0.1
}
\end{lstlisting}

Subsequent connections from a banned IP are rejected at layer 0 from the local ban
cache. No JA4 parsing occurs, hence the absent fingerprint fields.

\subsection{Config Reload Event}
\index{log format!config reload}

\begin{lstlisting}[language=json]
{
  "timestamp": "2026-04-04T13:00:00.000Z",
  "event": "config_reload",
  "action": null,
  "result": "success",
  "changed_keys": ["dial", "security_policy.spamhaus_bypass.enabled"],
  "duration_ms": 45.2
}
\end{lstlisting}

\subsection{Bypass Disabled Warning}
\index{log format!bypass disabled warning}

\begin{lstlisting}[language=json]
{
  "timestamp": "2026-04-04T08:00:00.001Z",
  "event": "bypass_disabled",
  "bypass": "alpn_browser_bypass",
  "effect": "browser traffic will be scored; false positive risk elevated",
  "level": "WARN"
}
\end{lstlisting}

\section{Human-Readable Text Format}
\index{log format!text format}

The text format is a condensed single-line summary derived from the JSON event.
It is written to a separate log stream (or can be toggled on/off) for operator
convenience during interactive sessions.

\begin{lstlisting}[language=bash]
ALLOWED:  172.19.0.10 | Country: IE  | JA4: t13d1113h2_... | Name: Browser (TLS 1.3) | TLS: TLS 1.3
BLOCKED:  185.220.101.1 | Country: DE | Reason: Spamhaus DROP: 185.220.96.0/21
TARPIT:   45.142.212.100 | Country: RU | JA4: t13d190900_... | Name: Sliver C2 | Score: 87
BANNED:   45.142.212.100 | Country: RU | Reason: Banned for 604800s (level 3)
\end{lstlisting}

\section{Log Levels}
\index{log format!log levels}

\begin{fullwidth}
\begin{table}[h]
\caption{Log levels and when each is emitted}
\label{tab:log-levels}
\begin{tabular}{@{}llp{8cm}@{}}
\toprule
\textbf{Level} & \textbf{Volume} & \textbf{Emitted for} \\
\midrule
ERROR & Low & Redis connection failure, config parse failure, TLS parser crash, external API repeated failure \\
WARN & Low & Bypass disabled at startup, AbuseIPDB rate limit hit, RDAP block expansion triggered \\
INFO & Medium & Config reload (success or error), startup completion, shutdown \\
DEBUG & High & Individual signal scores, cache hits/misses, Redis operations. Disabled in production. \\
(connection) & Very high & Every connection event — not strictly a log level; always emitted \\
\bottomrule
\end{tabular}
\end{table}
\end{fullwidth}

\section{Loki Log Queries (LogQL)}
\index{Loki}
\index{LogQL}

Assuming logs are shipped to Loki with the label \code{job="ja4proxy"}, the
following LogQL queries cover common operational needs:

\begin{lstlisting}[language=bash]
# All blocked connections in last hour
{job="ja4proxy"} | json | action="block"

# All connections from a specific IP
{job="ja4proxy"} | json | client_ip="185.220.101.1"

# Connections with risk_score > 80
{job="ja4proxy"} | json | risk_score > 80

# All Spamhaus bypasses
{job="ja4proxy"} | json | bypass="spamhaus"

# Tarpitted connections with their signal breakdown
{job="ja4proxy"} | json | action="tarpit"
  | line_format "{{.client_ip}} score={{.risk_score}}"

# Config reload events
{job="ja4proxy"} | json | event="config_reload"

# Error events only
{job="ja4proxy"} | json | level="ERROR"

# Fingerprint distribution for last 24h
{job="ja4proxy"} | json | fingerprint_name != ""
  | unwrap risk_score
  | rate by (fingerprint_name) [5m]
\end{lstlisting}

\section{SIEM Integration — CEF Format}
\index{SIEM integration}
\index{CEF format}

For SIEM platforms that require Common Event Format (CEF), the following mapping
converts \ja\ JSON fields to CEF:

\begin{fullwidth}
\begin{table}[h]
\caption{CEF field mapping for \ja\ connection events}
\label{tab:cef-mapping}
\begin{tabular}{@{}lll@{}}
\toprule
\textbf{CEF field} & \textbf{\ja\ JSON field} & \textbf{Notes} \\
\midrule
\code{CEF:0|JA4proxy|JA4proxy|1.0} & — & CEF header (vendor|product|version) \\
\code{name} & \code{event} + \code{action} & e.g.\ \code{connection\_block} \\
\code{severity} & \code{risk\_score} / 10 & Scale: 0–10 \\
\code{src} & \code{client\_ip} & Source IP \\
\code{cs1} (JA4) & \code{ja4} & Custom string field \\
\code{cs2} (Country) & \code{country} & Custom string field \\
\code{cs3} (Reason) & \code{reason} & Custom string field \\
\code{cn1} (Score) & \code{risk\_score} & Custom number field \\
\code{cn2} (Dial) & \code{dial} & Custom number field \\
\code{act} & \code{action} & Device action \\
\code{start} & \code{timestamp} & Event time \\
\code{deviceProcessName} & \code{"ja4proxy"} & Fixed string \\
\bottomrule
\end{tabular}
\end{table}
\end{fullwidth}

Example CEF line for a tarpitted connection:

\begin{lstlisting}[language=bash]
CEF:0|JA4proxy|JA4proxy|1.0|connection_tarpit|connection_tarpit|8|
  src=45.142.212.100 act=tarpit cs1Label=JA4
  cs1=t13d190900_9dc949149365_97f8aa674fd9
  cs2Label=Country cs2=RU cs3Label=Reason
  cs3=Score 87 exceeds tarpit threshold cn1Label=Score
  cn1=87 cn2Label=Dial cn2=75
  start=2026-04-04T12:36:42.100Z
\end{lstlisting}
